% Part: incompleteness
% Chapter: theories-computability
% Section: intepretability

\documentclass[../../../include/open-logic-section]{subfiles}

\begin{document}

\olfileid{inc}{tcp}{itp}

\olsection{$\Th{Q}$ دی تعبیر جنہاں نظرییاں وچ ہو سکدی اے، اوہ ناقابلِ فیصلہ نیں}

اسیں ایہناں نتیجیاں نوں ہور مضبوط بنا سکدے آں۔ غیر رسمی
طور اُتے اک زبان $\Lang{L_1}$ دی دوجی زبان $\Lang{L_2}$ وچ
تعبیر وچ $\Lang{L_1}$ دے دائرے، تعلق دیاں علامتاں تے فنکشن
دیاں علامتاں دی تعریف $\Lang{L_2}$ دے !!{formula}s نال کرنی
پیندی اے۔ اسیں ایتھے اوہ تفصیل نہیں دیواں گے، پر ایس تعریف
نوں پوری طرح ٹھیک ٹھیک بیان کیتا جا سکدا اے۔

\begin{thm}
فرض کرو $\Th{T}$ ایہو جیہی زبان دا نظریہ اے جیہدے وچ
حساب دی زبان دی تعبیر ہو سکدی اے، تے تعبیر ایہو جیہی اے کہ
$\Th{T}$، $\Th{Q}$ دی تعبیر نال کٹھا بے تضاد اے۔ تاں
$\Th{T}$ ناقابلِ فیصلہ اے۔ جے $\Th{T}$، $\Th{Q}$ دے
مسلّمیاں دیاں تعبیراں ثابت کردا اے، تاں $\Th{T}$ دی کوئی
بے تضاد توسیع قابلِ فیصلہ نہیں۔
\end{thm}

پچھلے قضیے دے ثبوت وچ تھوڑی جیہی تبدیلی نال ایہ ثبوت مل
جاندا اے۔ کسے مخالف مثال نال $\Th{Q}$ تے $\Th{\bar Q}$ نوں
قابلِ حساب طریقے نال وکھ کرن والا سیٹ مل جائے گا۔ $\Th{ZFC}$،
یعنی انتخاب دے مسلّمے والا زرمیلو--فرینکل سیٹ نظریہ، ایسی
مسلّمی بنیاد اے جیہڑی عام ریاضی دا خاصا حصہ کرن لئی کافی اے۔
$\Th{ZFC}$ وچ قدرتی اعداد دی تعریف ہو سکدی اے، تے ایس تعبیر
دے حساب نال $\Th{Q}$ دے مسلّمے سچے نیں۔ ایس توں ایہ نتیجہ
نکلدا اے۔

\begin{cor}
% PNB-CMP-069: اصل نتیجے وچ بے تضادی دی غائب شرط بحال؛اُتلے قضیے دی شرط مطابق۔
$\Th{ZFC}$ دی کوئی بے تضاد تے قابلِ فیصلہ توسیع نہیں۔
\end{cor}

\begin{cor}
$\Th{ZFC}$ دی کوئی توسیع ایہو جیہی نہیں جیہڑی مکمل، بے تضاد تے
قابلِ حساب طریقے نال !!{axiomatizable} ہووے۔
\end{cor}

$\Th{ZFC}$ دی زبان وچ اکّو دو رکنی تعلق دی علامت، $\in$، اے۔
(حقیقت وچ برابری دی وی لوڑ نہیں۔) ایس توں ایہ نتیجہ نکلدا اے۔

\begin{cor}
دو رکنی تعلق دی علامت والی کسے وی زبان دی پہلے درجے دی
منطق ناقابلِ فیصلہ اے۔
\end{cor}

ایہ نتیجہ دو اک رکنی فنکشناں دیاں علامتاں والی کسے وی زبان
لئی وی سچ اے، کیوں جے ایہناں دوہاں نال دو رکنی تعلق دی علامت
دی نقل کیتی جا سکدی اے۔ انہاں نتیجیاں دیاں حدّاں ٹھیک ایتھے
ای نیں۔ جے زبان وچ ساریاں تعلق دیاں علامتاں \emph{اک رکنی} ہون،
تے ودھ توں ودھ اک \emph{اک رکنی} فنکشن دی علامت ہووے، تاں
اوہدی پہلے درجے دی منطق قابلِ فیصلہ اے۔

اک ہور دلچسپ گل۔ اسیں جان دے آں کہ $\Obj 0$، $'$، $+$،
$\times$، $<$ والیاں علامتاں دی زبان دے اوہناں جملیاں دا سیٹ
ناقابلِ فیصلہ اے جیہڑے معیاری نمونے وچ سچے نیں۔ اصل وچ $<$
نوں باقی علامتاں نال تعریف کیتا جا سکدا اے، تے فیر $+$ نوں
$\times$ تے $'$ نال۔ ایس لئی $\Obj 0$، $'$، $\times$ والی
زبان دے سچے جملیاں دا سیٹ ناقابلِ فیصلہ اے۔ دوجے پاسے،
پریسبورگر نے وکھایا کہ $\Obj 0$، $'$، $+$ والی حسابی زبان
دے اوہناں جملیاں دا سیٹ قابلِ فیصلہ اے جیہڑے معیاری نمونے وچ
سچے نیں۔ پر ایہ فیصلہ کرن لئی اینا حساب کرنا پیندا اے کہ
طریقہ عملی طور اُتے ورتن جوگ نہیں۔

\end{document}
